\section{The Method of Frobenius II}

\noindent
In this section we discuss a method for finding two linearly
independent Frobenius solutions of a homogeneous linear second order
equation near a regular singular point in the case where the indicial
equation has a repeated real root. As in the preceding section, we
consider equations that can be written as
\begin{equation} \label{eq:7.6.1}
x^2(\alpha_0+\alpha_1x+\alpha_2x^2)y''+x(\beta_0+\beta_1x+\beta_2x^2)y'
+(\gamma_0+\gamma_1x+\gamma_2x^2)y=0
\end{equation}
where $\alpha_0\ne0$. We assume that the indicial equation $p_0(r)=0$
has a repeated real root $r_1$. In this case Theorem~\ref{thmtype:7.5.3}
implies that \eqref{eq:7.6.1} has one solution of the form
$$
y_1=x^{r_1}\sum_{n=0}^\infty a_nx^n,
$$
but does not provide a second solution $y_2$ such that $\{y_1,y_2\}$
is a fundamental set of solutions. The following extension of
Theorem~\ref{thmtype:7.5.2} provides a way to find a second solution.

\begin{theorem}\color{blue} \label{thmtype:7.6.1}
Let
\begin{equation} \label{eq:7.6.2}
Ly=
x^2(\alpha_0+\alpha_1x+\alpha_2x^2)y''+x(\beta_0+\beta_1x+\beta_2x^2)y'
+(\gamma_0+\gamma_1x+\gamma_2x^2)y,
\end{equation}
where $\alpha_0\ne0$ and define
\begin{eqnarray*}
p_0(r)&=&\alpha_0r(r-1)+\beta_0r+\gamma_0,\\[5pt]
p_1(r)&=&\alpha_1r(r-1)+\beta_1r+\gamma_1,\\[5pt]
p_2(r)&=&\alpha_2r(r-1)+\beta_2r+\gamma_2.\\
\end{eqnarray*}
Suppose $r$ is
a real number such that $p_0(n+r)$ is nonzero
for all positive integers $n$, and define
$$
\begin{array}{ccl}
a_0(r)&=&1,\\
a_1(r)&=&-\dst{p_1(r)\over p_0(r+1)},\\[10pt]
a_n(r)&=&-\dst{p_1(n+r-1)a_{n-1}(r)+p_2(n+r-2)a_{n-2}(r)\over
p_0(n+r)},\quad n\ge2.
\end{array}
$$
Then the Frobenius series
\begin{equation} \label{eq:7.6.3}
y(x,r)=x^r\sum_{n=0}^\infty a_n(r)x^n
\end{equation}
satisfies
\begin{equation} \label{eq:7.6.4}
Ly(x,r)=p_0(r)x^r.
\end{equation}
Moreover$,$
\begin{equation} \label{eq:7.6.5}
{\partial y\over \partial r}(x,r)=y(x,r)\ln x+x^r\sum_{n=1}^\infty
a_n'(r) x^n,
\end{equation}
and
\begin{equation} \label{eq:7.6.6}
L\left({\partial y\over \partial
r}(x,r)\right)=p'_0(r)x^r+x^rp_0(r)\ln x.
\end{equation}
\end{theorem}

\proof
Theorem~\ref{thmtype:7.5.2} implies \eqref{eq:7.6.4}. Differentiating
formally
with respect to $r$ in \eqref{eq:7.6.3} yields
\begin{eqnarray*}
{\partial y\over \partial
r}(x,r)&=&\dst{{\partial\over\partial r}(x^r)\sum_{n=0}^\infty
a_n(r)x^n +x^r\sum_{n=1}^\infty  a_n'(r)x^n}\\[10pt]
&=&\dst{x^r\ln x\sum_{n=0}^\infty
a_n(r)x^n +x^r\sum_{n=1}^\infty  a_n'(r)x^n}\\[10pt]
&=&y(x,r) \ln x +
x^r\sum_{n=1}^\infty  a_n'(r)x^n,
\end{eqnarray*}
which proves \eqref{eq:7.6.5}.

To prove that  $\partial y(x,r)/\partial r$ satisfies \eqref{eq:7.6.6}, we
view $y$ in \eqref{eq:7.6.2} as a function $y=y(x,r)$  of two
variables, where the prime indicates partial differentiation with
respect to $x$;   thus,
$$
y'=y'(x,r)={\partial y\over\partial x}(x,r)\mbox{\quad and \quad}
y''=y''(x,r)={\partial^2 y\over\partial x^2}(x,r).
$$
With this notation we can use \eqref{eq:7.6.2} to rewrite \eqref{eq:7.6.4} as
\begin{equation} \label{eq:7.6.7}
x^2q_0(x){\partial^2 y\over \partial x^2}(x,r)+xq_1(x){\partial y\over
\partial x}(x,r)+q_2(x)y(x,r)=p_0(r)x^r,
\end{equation}
where
\begin{eqnarray*}
q_0(x)&=&\alpha_0+\alpha_1x+\alpha_2x^2,\\[5pt]
q_1(x)&=&\beta_0+\beta_1x+\beta_2x^2,\\[5pt]
q_2(x)&=&\gamma_0+\gamma_1x+\gamma_2x^2.\\[5pt]
\end{eqnarray*}
Differentiating both sides of \eqref{eq:7.6.7} with respect to $r$
yields
$$
x^2q_0(x){\partial^3y\over \partial r\partial x^2}(x,r)+
xq_1(x){\partial^2y\over \partial r\partial x}(x,r)+q_2(x){\partial
y\over\partial r}(x,r)=p'_0(r)x^r+p_0(r) x^r \ln x.
$$
By changing the order of differentiation in the first two terms on the left
we can rewrite this  as
$$
x^2q_0(x){\partial^3 y\over \partial x^2\partial r}(x,r)
+xq_1(x){\partial^2 y\over \partial x\partial r}(x,r)+q_2(x){\partial
y\over \partial r}(x,r)=p'_0(r)x^r+p_0(r) x^r \ln x,
$$
or
$$
x^2q_0(x){\partial^2\over \partial x^2}
\left({\partial y\over\partial r}(x,r)\right)
+xq_1(x){\partial\over\partial r}\left({\partial y\over\partial
x}(x,r)\right) +q_2(x){\partial y\over\partial r}(x,r)=
p'_0(r)x^r+p_0(r) x^r \ln x,
$$
which is equivalent to \eqref{eq:7.6.6}.

\begin{theorem}\color{blue} \label{thmtype:7.6.2}
Let $L$ be as in Theorem~{\rm\ref{thmtype:7.6.1}} and suppose the
indicial equation $p_0(r)=0$ has a repeated real root $r_1.$
Then
$$
y_1(x)=y(x,r_1)=x^{r_1}\sum_{n=0}^\infty a_n(r_1)x^n
$$
and
\begin{equation} \label{eq:7.6.8}
y_2(x)={\partial y\over\partial r}(x,r_1)=y_1(x)\ln
x+x^{r_1}\sum_{n=1}^\infty a_n'(r_1)x^n
\end{equation}
form a fundamental set of solutions of $Ly=0.$
\end{theorem}

\proof
Since $r_1$  is a repeated root of $p_0(r)=0$, the indicial polynomial
can be factored as
$$
p_0(r)=\alpha_0(r-r_1)^2,
$$
so
$$
p_0(n+r_1)=\alpha_0n^2,
$$
which is nonzero if $n>0$. Therefore the assumptions of
Theorem~\ref{thmtype:7.6.1} hold with $r=r_1$, and \eqref{eq:7.6.4} implies
that $Ly_1=p_0(r_1)x^{r_1}=0$. Since
$$
p_0'(r)=2\alpha(r-r_1)
$$
it follows that $p_0'(r_1)=0$, so \eqref{eq:7.6.6} implies that
$$
Ly_2=p_0'(r_1)x^{r_1}+x^{r_1}p_0(r_1)\ln x=0.
$$
This proves that $y_1$ and $y_2$ are both solutions of $Ly=0$. We
leave the proof that $\{y_1,y_2\}$ is a fundamental set as an exercise
(Exercise~\ref{exer:7.6.53}).

\begin{example}\label{example:7.6.1} \rm
Find  a fundamental set of solutions of
\begin{equation} \label{eq:7.6.9}
x^2(1-2x+x^2)y''-x(3+x)y'+(4+x)y=0.
\end{equation}
Compute just the terms involving $x^{n+r_1}$, where  $0\le n\le4$
and $r_1$ is the root of the indicial equation.
\end{example}

\solution
For the given equation,  the polynomials defined in
Theorem~\ref{thmtype:7.6.1} are
$$
\begin{array}{lllll}
p_0(r)&=&r(r-1)-3r+4&=&(r-2)^2,\\[5pt]
p_1(r)&=&-2r(r-1)-r+1&=&-(r-1)(2r+1),\\[5pt]
p_2(r)&=&r(r-1).
\end{array}
$$
Since $r_1=2$ is a repeated root of the indicial polynomial $p_0$,
Theorem~\ref{thmtype:7.6.2} implies that
\begin{equation} \label{eq:7.6.10}
y_1=x^2\sum_{n=0}^\infty  a_n(2)x^n\quad\mbox{ and }\quad
y_2=y_1\ln x+x^2\sum_{n=1}^\infty a_n'(2)x^n
\end{equation}
form a fundamental set of  Frobenius solutions of \eqref{eq:7.6.9}.
To find the coefficients in these series, we use the recurrence
formulas from Theorem~\ref{thmtype:7.6.1}:
\begin{equation} \label{eq:7.6.11}
\begin{array}{ccl}
a_0(r)&=&1,\\
a_1(r)&=&-\dst{p_1(r)\over p_0(r+1)}
=-\dst{(r-1)(2r+1)\over(r-1)^2}
=\dst{2r+1\over r-1},\\[10pt]
a_n(r)&=&-\dst{p_1(n+r-1)a_{n-1}(r)+p_2(n+r-2)a_{n-2}(r)\over
p_0(n+r)}\\[10pt]
&=&\dst{(n+r-2)\left[(2n+2r-1)a_{n-1}(r)
-(n+r-3)a_{n-2}(r)\right]\over(n+r-2)^2}\\[10pt]
&=&\dst{{(2n+2r-1)\over(n+r-2)}a_{n-1}(r)-
{(n+r-3)\over(n+r-2)}a_{n-2}(r)},n\ge2.
\end{array}
\end{equation}
Differentiating yields
\begin{equation} \label{eq:7.6.12}
\begin{array}{ccl}
a'_1(r)&=&-\dst{3\over (r-1)^2},\\[10pt]
a'_n(r)&=&\dst{{2n+2r-1\over n+r-2}a'_{n-1}(r)-{n+r-3\over
n+r-2}a'_{n-2}(r)}\\[10pt]
&&\dst{-{3\over(n+r-2)^2}a_{n-1}(r)-{1\over(n+r-2)^2}a_{n-2}(r)},\quad
n\ge2.
\end{array}
\end{equation}

Setting $r=2$ in \eqref{eq:7.6.11} and \eqref{eq:7.6.12} yields
\begin{equation} \label{eq:7.6.13}
\begin{array}{ccl}
a_0(2)&=&1,\\
a_1(2)&=&5,\\[10pt]
a_n(2)&=&\dst{{(2n+3)\over n}
a_{n-1}(2)-{(n-1)\over n}a_{n-2}(2)},\quad n\ge2
\end{array}
\end{equation}
and
\begin{equation} \label{eq:7.6.14}
\begin{array}{ccl}
a_1'(2)&=&-3,\\[10pt]
a'_n(2)&=&\dst{{2n+3\over n}a'_{n-1}(2)-{n-1\over n}a'_{n-2}(2)
-{3\over n^2}a_{n-1}(2)-{1\over n^2}a_{n-2}(2)},\quad n\ge2.
\end{array}
\end{equation}
Computing recursively with  \eqref{eq:7.6.13} and \eqref{eq:7.6.14}
yields
$$
a_0(2)=1,\,a_1(2)=5,\,a_2(2)=17,\,a_3(2)={143\over3},\,a_4(2)={355\over3},
$$
and
$$
a_1'(2)=-3,\,a_2'(2)=-{29\over2},\,a_3'(2)=-{859\over18},
\,a_4'(2)=-{4693\over36}.
$$
Substituting these coefficients into \eqref{eq:7.6.10} yields
$$
y_1=x^2\left(1+5x+17x^2+{143\over3}x^3
+{355\over3}x^4+\cdots\right)
$$
and
$$
y_2=y_1 \ln x
-x^3\left(3+{29\over2}x+{859\over18}x^2+{4693\over36}x^3
+\cdots\right). \bbox
$$

Since the recurrence formula \eqref{eq:7.6.11} involves three terms, it's
not possible to obtain a simple explicit formula for the coefficients
in the Frobenius solutions of \eqref{eq:7.6.9}. However, as we saw in the
preceding sections, the recurrrence formula for $\{a_n(r)\}$ involves
only two terms if either $\alpha_1=\beta_1=\gamma_1=0$ or
$\alpha_2=\beta_2=\gamma_2=0$ in \eqref{eq:7.6.1}. In this case, it's
often
possible to find explicit formulas for the coefficients. The next two
examples illustrate this.

\begin{example}\label{example:7.6.2} \rm
Find a fundamental set of Frobenius  solutions of
\begin{equation} \label{eq:7.6.15}
2x^2(2+x)y''+5x^2y'+(1+x)y=0.
\end{equation}
Give explicit formulas for the coefficients in the solutions.
\end{example}

\solution
For  the given equation, the polynomials defined in
Theorem~\ref{thmtype:7.6.1} are
$$
\begin{array}{ccccc}
p_0(r)&=&4r(r-1)+1&=&(2r-1)^2,\\[5pt]
p_1(r)&=&2r(r-1)+5r+1&=&(r+1)(2r+1),\\[5pt]
p_2(r)&=&0.
\end{array}
$$
Since $r_1=1/2$ is a repeated zero of the  indicial polynomial $p_0$,
Theorem~\ref{thmtype:7.6.2} implies that
\begin{equation} \label{eq:7.6.16}
y_1=x^{1/2}\sum_{n=0}^\infty  a_n(1/2)x^n
\end{equation}
and
\begin{equation} \label{eq:7.6.17}
y_2=y_1\ln x+x^{1/2}\sum_{n=1}^\infty a_n'(1/2)x^n
\end{equation}
form a fundamental set of Frobenius solutions of \eqref{eq:7.6.15}.
Since $p_2\equiv0$, the recurrence formulas in Theorem~\ref{thmtype:7.6.1}
 reduce to
$$
\begin{array}{ccl}
a_0(r)&=&1,\\
a_n(r)&=&-\dst{p_1(n+r-1)\over p_0(n+r)}a_{n-1}(r),\\[10pt]
&=&-\dst{(n+r)(2n+2r-1)\over(2n+2r-1)^2}a_{n-1}(r),\\[10pt]
&=&-\dst{n+r\over2n+2r-1}a_{n-1}(r),\quad n\ge0.
\end{array}
$$
We leave it to you to show that
\begin{equation} \label{eq:7.6.18}
a_n(r)=(-1)^n\prod_{j=1}^n{j+r\over2j+2r-1},\quad n\ge0.
\end{equation}
Setting $r=1/2$  yields
\begin{equation} \label{eq:7.6.19}
\begin{array}{ccl}
a_n(1/2)&=&(-1)^n\dst\prod_{j=1}^n\dst{j+1/2\over2j}=
(-1)^n\prod_{j=1}^n\dst{2j+1\over4j},\\[10pt]
&=&\dst{(-1)^n\prod_{j=1}^n(2j+1)\over4^nn!},\quad n\ge0.
\end{array}
\end{equation}
Substituting this into \eqref{eq:7.6.16} yields
$$
y_1=x^{1/2}\sum_{n=0}^\infty{(-1)^n\prod_{j=1}^n(2j+1)\over4^nn!}x^n.
$$

To obtain $y_2$ in \eqref{eq:7.6.17}, we must compute $a_n'(1/2)$
for $n=1$, $2$,\dots. We'll do this by logarithmic
differentiation.  From \eqref{eq:7.6.18},
$$
|a_n(r)|=\prod_{j=1}^n{|j+r|\over|2j+2r-1|},\quad n\ge1.
$$
Therefore
$$
\ln |a_n(r)|=\sum^n_{j=1} \left(\ln |j+r|-\ln|2j+2r-1|\right).
$$
Differentiating with respect to $r$ yields
$$
{a'_n(r)\over a_n(r)}=\sum^n_{j=1} \left({1\over
j+r}-{2\over2j+2r-1}\right).
$$
Therefore
$$
a'_n(r)=a_n(r) \sum^n_{j=1} \left({1\over j+r}-{2\over2j+2r-1}\right).
$$
Setting $r=1/2$ here and recalling \eqref{eq:7.6.19} yields
\begin{equation} \label{eq:7.6.20}
a'_n(1/2)={(-1)^n\prod_{j=1}^n(2j+1)\over4^nn!}\left(\sum_{j=1}^n{1\over
j+1/2}-\sum_{j=1}^n{1\over j}\right).
\end{equation}
Since
$$
{1\over j+1/2}-{1\over j}={j-j-1/2\over j(j+1/2)}=-{1\over j(2j+1)},
$$
 \eqref{eq:7.6.20} can be rewritten as
$$
a'_n(1/2)=-{(-1)^n\prod_{j=1}^n(2j+1)\over4^nn!}
\sum_{j=1}^n{1\over j(2j+1)}.
$$
Therefore, from \eqref{eq:7.6.17},
$$
y_2=y_1\ln
x-x^{1/2}\sum_{n=1}^\infty{(-1)^n\prod_{j=1}^n(2j+1)\over4^nn!}
\left(\sum_{j=1}^n{1\over j(2j+1)}\right)x^n.
$$

\begin{example}\label{example:7.6.3} \rm
Find a fundamental set of Frobenius  solutions of
\begin{equation} \label{eq:7.6.21}
x^2(2-x^2)y''-2x(1+2x^2)y'+(2-2x^2)y=0.
\end{equation}
Give explicit formulas for the coefficients in the solutions.
\end{example}

\solution
For \eqref{eq:7.6.21}, the polynomials defined in
Theorem~\ref{thmtype:7.6.1} are
$$
\begin{array}{ccccc}
p_0(r)&=&2r(r-1)-2r+2&=&2(r-1)^2,\\[5pt]
p_1(r)&=&0,\\[5pt]
p_2(r)&=&-r(r-1)-4r-2&=&-(r+1)(r+2).
\end{array}
$$
As in Section~7.5, since $p_1\equiv0$,  the recurrence
formulas
of Theorem~\ref{thmtype:7.6.1} imply that $a_n(r)=0$ if $n$ is odd, and
$$
\begin{array}{ccl}
a_0(r)&=&1,\\
a_{2m}(r)&=&-\dst{p_2(2m+r-2)\over p_0(2m+r)}a_{2m-2}(r)\\[10pt]
&=&\dst{(2m+r-1)(2m+r)\over2(2m+r-1)^2}a_{2m-2}(r)\\[10pt]
&=&\dst{2m+r\over2(2m+r-1)}a_{2m-2}(r),\quad m\ge1.
\end{array}
$$
Since $r_1=1$ is a repeated root of the indicial polynomial $p_0$,
Theorem~\ref{thmtype:7.6.2} implies that
\begin{equation} \label{eq:7.6.22}
y_1=x\sum_{m=0}^\infty  a_{2m}(1)x^{2m}
\end{equation}
and
\begin{equation} \label{eq:7.6.23}
y_2=y_1\ln x+x\sum_{m=1}^\infty a'_{2m}(1)x^{2m}
\end{equation}
form a fundamental set of  Frobenius solutions of \eqref{eq:7.6.21}.
We leave it to you to show  that
\begin{equation} \label{eq:7.6.24}
a_{2m}(r)={1\over2^m}\prod_{j=1}^m{2j+r\over2j+r-1}.
\end{equation}
Setting $r=1$ yields
\begin{equation} \label{eq:7.6.25}
a_{2m}(1)={1\over2^m}\prod_{j=1}^m{2j+1\over2j}
={\prod_{j=1}^m(2j+1)\over4^mm!},
\end{equation}
and substituting this into \eqref{eq:7.6.22} yields
$$
y_1=x\sum_{m=0}^\infty{\prod_{j=1}^m(2j+1)\over4^mm!}x^{2m}.
$$

To obtain $y_2$ in \eqref{eq:7.6.23}, we must compute $a_{2m}'(1)$
for $m=1$, $2$, \dots. Again we use logarithmic differentiation. From
\eqref{eq:7.6.24},
$$
|a_{2m}(r)|={1\over2^m}\prod_{j=1}^m{|2j+r|\over|2j+r-1|}.
$$
Taking logarithms yields
$$
\ln |a_{2m}(r)|=-m\ln2+ \sum^m_{j=1} \left(\ln
|2j+r|-\ln|2j+r-1|\right).
$$
Differentiating with respect to $r$ yields
$$
{a'_{2m}(r)\over a_{2m}(r)}=\sum^m_{j=1} \left({1\over
2j+r}-{1\over2j+r-1}\right).
$$
Therefore
$$
a'_{2m}(r)=a_{2m}(r) \sum^m_{j=1} \left({1\over
2j+r}-{1\over2j+r-1}\right).
$$
Setting $r=1$  and recalling \eqref{eq:7.6.25} yields
\begin{equation} \label{eq:7.6.26}
a'_{2m}(1)=\dst{{\prod_{j=1}^m(2j+1)\over4^mm!}
\sum_{j=1}^m\left({1\over2j+1}-{1\over2j}\right)}.
\end{equation}
Since
$$
{1\over2j+1}-{1\over2j}=-{1\over2j(2j+1)},
$$
\eqref{eq:7.6.26} can be rewritten as
$$
a_{2m}'(1)=-\dst{{\prod_{j=1}^m(2j+1)\over2\cdot4^mm!}
\sum_{j=1}^m{1\over j(2j+1)}}.
$$
Substituting this into \eqref{eq:7.6.23} yields
$$
y_2=y_1\ln
x-{x\over2}\sum_{m=1}^\infty{\prod_{j=1}^m(2j+1)\over4^mm!}
\left(\sum_{j=1}^m{1\over j(2j+1)}\right)x^{2m}. \bbox
$$

If the solution $y_1=y(x,r_1)$ of $Ly=0$ reduces to a finite sum, then
there's a
difficulty in using logarithmic differentiation to obtain the coefficients
$\{a_n'(r_1)\}$ in the second solution. The next example illustrates
this difficulty and shows how to overcome it.

\begin{example}\label{example:7.6.4} \rm
Find a fundamental set of Frobenius solutions of
\begin{equation} \label{eq:7.6.27}
x^2y''-x(5-x)y'+(9-4x)y=0.
\end{equation}
Give explicit formulas for the coefficients in the solutions.
\end{example}

\solution
For \eqref{eq:7.6.27} the polynomials defined in Theorem~\ref{thmtype:7.6.1}
are
$$
\begin{array}{ccccc}
p_0(r)&=&r(r-1)-5r+9&=&(r-3)^2,\\[5pt]
p_1(r)&=&r-4,\\[5pt]
p_2(r)&=&0.
\end{array}
$$
Since $r_1=3$ is a repeated zero of the  indicial polynomial $p_0$,
Theorem~\ref{thmtype:7.6.2} implies that
\begin{equation} \label{eq:7.6.28}
y_1=x^3\sum_{n=0}^\infty  a_n(3)x^n
\end{equation}
and
\begin{equation} \label{eq:7.6.29}
y_2=y_1\ln x+x^3\sum_{n=1}^\infty a_n'(3)x^n
\end{equation}
are linearly independent Frobenius solutions of \eqref{eq:7.6.27}.
To find the coefficients in \eqref{eq:7.6.28} we use the recurrence
formulas
$$
\begin{array}{ccl}
a_0(r)&=&1,\\
a_n(r)&=&-\dst{p_1(n+r-1)\over p_0(n+r)}a_{n-1}(r)\\[10pt]
&=&-\dst{n+r-5\over(n+r-3)^2}a_{n-1}(r),\quad n\ge1.
\end{array}
$$
We leave it to you to show that
\begin{equation} \label{eq:7.6.30}
a_n(r)=(-1)^n\prod_{j=1}^n{j+r-5\over(j+r-3)^2}.
\end{equation}
Setting $r=3$ here yields
$$
a_n(3)=(-1)^n\prod_{j=1}^n{j-2\over j^2},
$$
so  $a_1(3)=1$ and $a_n(3)=0$ if $n\ge2$. Substituting
these coefficients into \eqref{eq:7.6.28} yields
$$
y_1=x^3(1+x).
$$

To obtain $y_2$ in \eqref{eq:7.6.29} we must compute $a_n'(3)$
for $n=1$, $2$, \dots.
 Let's first try  logarithmic differentiation.
From \eqref{eq:7.6.30},
$$
|a_n(r)|=\prod_{j=1}^n{|j+r-5|\over|j+r-3|^2},\quad n\ge1,
$$
so
$$
\ln |a_n(r)|=\sum^n_{j=1} \left(\ln |j+r-5|-2\ln|j+r-3|\right).
$$
Differentiating  with respect to $r$ yields
$$
{a'_n(r)\over a_n(r)}=\sum^n_{j=1} \left({1\over
j+r-5}-{2\over j+r-3}\right).
$$
Therefore
\begin{equation} \label{eq:7.6.31}
a'_n(r)=a_n(r) \sum^n_{j=1} \left({1\over j+r-5}-{2\over j+r-3}\right).
\end{equation}
However, we can't simply set $r=3$ here if $n\ge2$,  since the bracketed
expression in the sum corresponding to  $j=2$ contains the term $1/(r-3)$.
In fact, since $a_n(3)=0$ for $n\ge2$, the  formula \eqref{eq:7.6.31} for
$a_n'(r)$ is actually an indeterminate form at $r=3$.

We  overcome this difficulty as follows. From \eqref{eq:7.6.30} with
$n=1$,
$$
a_1(r)=-{r-4\over (r-2)^2}.
$$
Therefore
$$
a_1'(r)={r-6\over(r-2)^3},
$$
so
\begin{equation} \label{eq:7.6.32}
a_1'(3)=-3.
\end{equation}
From \eqref{eq:7.6.30} with $n\ge2$,
$$
a_n(r)=(-1)^n (r-4)(r-3)\,{\prod_{j=3}^n(j+r-5)\over\prod_{j=1}^n(j+r-3)^2}
=(r-3)c_n(r),
$$
where
$$
c_n(r)=(-1)^n(r-4)\,
{\prod_{j=3}^n(j+r-5)\over\prod_{j=1}^n(j+r-3)^2},\quad  n\ge2.
$$
Therefore
$$
a_n'(r)=c_n(r)+(r-3)c_n'(r),\quad n\ge2,
$$
which implies that  $a_n'(3)=c_n(3)$ if $n\ge3$. We leave it to
you to verify that
$$
a_n'(3)=c_n(3)={(-1)^{n+1}\over n(n-1)n!},\quad n\ge2.
$$
Substituting this and \eqref{eq:7.6.32} into \eqref{eq:7.6.29}
yields
$$
y_2=x^3(1+x)\ln x-3x^4-x^3\dst{\sum_{n=2}^\infty {(-1)^n\over
n(n-1)n!}x^n}.
$$

\exercises
In Exercises~\ref{exer:7.6.1}--\ref{exer:7.6.11} find a fundamental set of
Frobenius solutions.
Compute  the terms involving $x^{n+r_1}$,
where $0\le n\le N$ ($N$ at least $7$) and $r_1$ is the root of the
indicial equation.
 Optionally, write a computer program to implement the applicable
recurrence formulas and take $N>7$.

\begin{exerciselist}

\item\label{exer:7.6.1} \Cex
$x^2y''-x(1-x)y'+(1-x^2)y=0$

\item\label{exer:7.6.2} \Cex
$x^2(1+x+2x^2)y'+x(3+6x+7x^2)y'+(1+6x-3x^2)y=0$

\item\label{exer:7.6.3} \Cex
$x^2(1+2x+x^2)y''+x(1+3x+4x^2)y'-x(1-2x)y=0$

\item\label{exer:7.6.4} \Cex
$4x^2(1+x+x^2)y''+12x^2(1+x)y'+(1+3x+3x^2)y=0$

\item\label{exer:7.6.5} \Cex
$x^2(1+x+x^2)y''-x(1-4x-2x^2)y'+y=0$

\item\label{exer:7.6.6} \Cex
$9x^2y''+3x(5+3x-2x^2)y'+(1+12x-14x^2)y=0$

\item\label{exer:7.6.7} \Cex
$x^2y''+x(1+x+x^2)y'+x(2-x)y=0$

\item\label{exer:7.6.8} \Cex
$x^2(1+2x)y''+x(5+14x+3x^2)y'+(4+18x+12x^2)y=0$

\item\label{exer:7.6.9} \Cex
$4x^2y''+2x(4+x+x^2)y'+(1+5x+3x^2)y=0$

\item\label{exer:7.6.10} \Cex
$16x^2y''+4x(6+x+2x^2)y'+(1+5x+18x^2)y=0$

\item\label{exer:7.6.11} \Cex
$9x^2(1+x)y''+3x(5+11x-x^2)y'+(1+16x-7x^2)y=0$

\exercisetext{In Exercises~\ref{exer:7.6.12}--\ref{exer:7.6.22}
find a fundamental set of Frobenius  solutions.
Give explicit formulas for the coefficients.}

\item\label{exer:7.6.12}
$4x^2y''+(1+4x)y=0$

\item\label{exer:7.6.13}
$36x^2(1-2x)y''+24x(1-9x)y'+(1-70x)y=0$

\item\label{exer:7.6.14}
$x^2(1+x)y''-x(3-x)y'+4y=0$

\item\label{exer:7.6.15}
$x^2(1-2x)y''-x(5-4x)y'+(9-4x)y=0$

\item\label{exer:7.6.16}
$25x^2y''+x(15+x)y'+(1+x)y=0$

\item\label{exer:7.6.17}
$2x^2(2+x)y''+x^2y'+(1-x)y=0$

\item\label{exer:7.6.18}
$x^2(9+4x)y''+3xy'+(1+x)y=0$

\item\label{exer:7.6.19}
$x^2y''-x(3-2x)y'+(4+3x)y=0$

\item\label{exer:7.6.20}
$x^2(1-4x)y''+3x(1-6x)y'+(1-12x)y=0$

\item\label{exer:7.6.21}
$x^2(1+2x)y''+x(3+5x)y'+(1-2x)y=0$

\item\label{exer:7.6.22}
$2x^2(1+x)y''-x(6-x)y'+(8-x)y=0$

\exercisetext{In Exercises~\ref{exer:7.6.23}--\ref{exer:7.6.27} find a
fundamental set of Frobenius solutions. Compute the terms involving
$x^{n+r_1}$, where $0\le n\le N$ ($N$ at least $7$) and $r_1$ is the
root of the indicial equation. Optionally, write a computer program to
implement the applicable recurrence formulas and take $N>7$.}

\item\label{exer:7.6.23} \Cex $x^2(1+2x)y''+x(5+9x)y'+(4+3x)y=0$

\item\label{exer:7.6.24} \Cex $x^2(1-2x)y''-x(5+4x)y'+(9+4x)y=0$

\item\label{exer:7.6.25} \Cex
$x^2(1+4x)y''-x(1-4x)y'+(1+x)y=0$

\item\label{exer:7.6.26} \Cex
$x^2(1+x)y''+x(1+2x)y'+xy=0$

\item\label{exer:7.6.27} \Cex
$x^2(1-x)y''+x(7+x)y'+(9-x)y=0$

\exercisetext{In Exercises~\ref{exer:7.6.28}--\ref{exer:7.6.38} find a
fundamental set of Frobenius solutions. Give explicit formulas for the
coefficients.}

\item\label{exer:7.6.28}
$x^2y''-x(1-x^2)y'+(1+x^2)y=0$

\item\label{exer:7.6.29}
$x^2(1+x^2)y''-3x(1-x^2)y'+4y=0$

\item\label{exer:7.6.30}
$4x^2y''+2x^3y'+(1+3x^2)y=0$

\item\label{exer:7.6.31}
$x^2(1+x^2)y''-x(1-2x^2)y'+y=0$

\item\label{exer:7.6.32}
$2x^2(2+x^2)y''+7x^3y'+(1+3x^2)y=0$

\item\label{exer:7.6.33}
$x^2(1+x^2)y''-x(1-4x^2)y'+(1+2x^2)y=0$

\item\label{exer:7.6.34}
$4x^2(4+x^2)y''+3x(8+3x^2)y'+(1-9x^2)y=0$

\item\label{exer:7.6.35}
$3x^2(3+x^2)y''+x(3+11x^2)y'+(1+5x^2)y=0$

\item\label{exer:7.6.36}
$4x^2(1+4x^2)y''+32x^3y'+y=0$

\item\label{exer:7.6.37}
$9x^2y''-3x(7-2x^2)y'+(25+2x^2)y=0$

\item\label{exer:7.6.38}
$x^2(1+2x^2)y''+x(3+7x^2)y'+(1-3x^2)y=0$

\exercisetext{In Exercises~\ref{exer:7.6.39}--\ref{exer:7.6.43} find a
fundamental set of Frobenius solutions. Compute the terms involving
$x^{2m+r_1}$, where $0\le m\le M$ ($M$ at least $3$) and $r_1$ is the
root of the indicial equation. Optionally, write a computer program to
implement the applicable recurrence formulas and take $M>3$.}

\item\label{exer:7.6.39} \Cex
$x^2(1+x^2)y''+x(3+8x^2)y'+(1+12x^2)y$

\item\label{exer:7.6.40} \Cex
$x^2y''-x(1-x^2)y'+(1+x^2)y=0$

\item\label{exer:7.6.41} \Cex
$x^2(1-2x^2)y''+x(5-9x^2)y'+(4-3x^2)y=0$

\item\label{exer:7.6.42} \Cex
$x^2(2+x^2)y''+x(14-x^2)y'+2(9+x^2)y=0$

\item\label{exer:7.6.43} \Cex
$x^2(1+x^2)y''+x(3+7x^2)y'+(1+8x^2)y=0$

\exercisetext{In Exercises~\ref{exer:7.6.44}--\ref{exer:7.6.52} find a
fundamental set of Frobenius solutions. Give explicit formulas for the
coefficients.}

\item\label{exer:7.6.44}
$x^2(1-2x)y''+3xy'+(1+4x)y=0$

\item\label{exer:7.6.45}
$x(1+x)y''+(1-x)y'+y=0$

\item\label{exer:7.6.46}
$x^2(1-x)y''+x(3-2x)y'+(1+2x)y=0$

\item\label{exer:7.6.47}
$4x^2(1+x)y''-4x^2y'+(1-5x)y=0$

\item\label{exer:7.6.48}
$x^2(1-x)y''-x(3-5x)y'+(4-5x)y=0$

\item\label{exer:7.6.49}
$x^2(1+x^2)y''-x(1+9x^2)y'+(1+25x^2)y=0$

\item\label{exer:7.6.50}
$9x^2y''+3x(1-x^2)y'+(1+7x^2)y=0$

\item\label{exer:7.6.51}
$x(1+x^2)y''+(1-x^2)y'-8xy=0$

\item\label{exer:7.6.52}
$4x^2y''+2x(4-x^2)y'+(1+7x^2)y=0$

\item\label{exer:7.6.53}
Under the assumptions of Theorem~\ref{thmtype:7.6.2}, suppose
the power series
$$
\sum_{n=0}^\infty a_n(r_1)x^n \quad\mbox{ and }\quad
\sum_{n=1}^\infty a_n'(r_1)x^n
$$
  converge on $(-\rho,\rho)$.
\begin{alist}
\item % (a)
Show that
$$
y_1=x^{r_1}\sum_{n=0}^\infty a_n(r_1)x^n\quad\mbox{ and }\quad
y_2=y_1\ln x+x^{r_1}\sum_{n=1}^\infty a_n'(r_1)x^n
$$
are linearly independent on $(0,\rho)$. \hint{Show that if $c_1$
and $c_2$ are constants such that $c_1y_1+c_2y_2\equiv0$ on $(0,\rho)$,
then
$$
(c_1+c_2\ln x)\sum_{n=0}^\infty a_n(r_1)x^n+
c_2\sum_{n=1}^\infty a_n'(r_1)x^n=0,\quad 0<x<\rho.
$$
Then let $x\to0+$ to conclude that $c_2=0$.}
\item % (b)
Use the result of \part{a} to
complete the proof of Theorem~\ref{thmtype:7.6.2}.
\end{alist}

\item\label{exer:7.6.54}
Let
$$
Ly=x^2(\alpha_0+\alpha_1x)y''+x(\beta_0+\beta_1x)y'+(\gamma_0+\gamma_1x)y
$$
and define
$$
p_0(r)=\alpha_0r(r-1)+\beta_0r+\gamma_0\quad\mbox{ and }\quad
p_1(r)=\alpha_1r(r-1)+\beta_1r+\gamma_1.
$$
Theorem~\ref{thmtype:7.6.1}  and Exercise~7.5.~\hspace*{-3pt}\ref{exer:7.5.55}\part{a}
imply that if
$$
y(x,r)=x^r\sum_{n=0}^\infty a_n(r)x^n
$$
where
$$
a_n(r)=(-1)^n\prod_{j=1}^n{p_1(j+r-1)\over p_0(j+r)},
$$
then
$$
Ly(x,r)=p_0(r)x^r.
$$
Now suppose
$p_0(r)=\alpha_0(r-r_1)^2$ and $p_1(k+r_1)\ne0$ if $k$
is a nonnegative integer.
\begin{alist}
\item % (a)
Show that $Ly=0$ has the solution
$$
y_1=x^{r_1}\sum_{n=0}^\infty a_n(r_1)x^n,
$$
where
$$
a_n(r_1)={(-1)^n\over\alpha_0^n(n!)^2}\prod_{j=1}^np_1(j+r_1-1).
$$
\item % ()
Show that $Ly=0$ has the second solution
$$
y_2=y_1\ln x+x^{r_1}\sum_{n=1}^\infty a_n(r_1)J_nx^n,
$$
where
$$
J_n=\sum_{j=1}^n{p_1'(j+r_1-1)\over p_1(j+r_1-1)}-2\sum_{j=1}^n{1\over j}.
$$
\item % (c)
Conclude from \part{a} and \part{b} that if
 $\gamma_1\ne0$ then
$$
y_1=x^{r_1}\sum_{n=0}^\infty
{(-1)^n\over(n!)^2}\left(\gamma_1\over\alpha_0\right)^nx^n
$$
and
$$
y_2=y_1\ln x-2x^{r_1}\sum_{n=1}^\infty
{(-1)^n\over(n!)^2}\left(\gamma_1\over\alpha_0\right)^n
\left(\sum_{j=1}^n{1\over j}\right)x^n
$$
are solutions of
$$
\alpha_0x^2y''+\beta_0xy'+(\gamma_0+\gamma_1x)y=0.
$$
(The conclusion is also valid if $\gamma_1=0$. Why?)
\end{alist}

\item\label{exer:7.6.55}
Let
$$
Ly=x^2(\alpha_0+\alpha_qx^q)y''+x(\beta_0+\beta_qx^q)y'+(\gamma_0+\gamma_qx^q)y
$$
where $q$ is a positive integer, and define
$$
p_0(r)=\alpha_0r(r-1)+\beta_0r+\gamma_0\quad\mbox{ and }\quad
p_q(r)=\alpha_qr(r-1)+\beta_qr+\gamma_q.
$$
Suppose
$$
p_0(r)=\alpha_0(r-r_1)^2 \quad\mbox{ and }\quad p_q(r)\not\equiv0.
$$
\begin{alist}
\item % (a)
Recall from  Exercise~7.5.~\hspace*{-3pt}\ref{exer:7.5.59}  that
 $Ly~=0$ has the solution
$$
y_1=x^{r_1}\sum_{m=0}^\infty a_{qm}(r_1)x^{qm},
$$
where
$$
a_{qm}(r_1)={(-1)^m\over
(q^2\alpha_0)^m(m!)^2}\prod_{j=1}^mp_q\left(q(j-1)+r_1\right).
$$
\item % ()
Show that $Ly=0$ has the second solution
$$
y_2=y_1\ln x+x^{r_1}\sum_{m=1}^\infty a_{qm}'(r_1)J_mx^{qm},
$$
where
$$
J_m=\sum_{j=1}^m{p_q'\left(q(j-1)+r_1\right)\over
p_q\left(q(j-1)+r_1\right)}-{2\over q}\sum_{j=1}^m{1\over j}.
$$
\item % (c)
Conclude from \part{a} and \part{b} that if
 $\gamma_q\ne0$ then
$$
y_1=x^{r_1}\sum_{m=0}^\infty
{(-1)^m\over(m!)^2}\left(\gamma_q\over q^2\alpha_0\right)^mx^{qm}
$$
and
$$
y_2=y_1\ln x-{2\over q}x^{r_1}\sum_{m=1}^\infty
{(-1)^m\over(m!)^2}\left(\gamma_q\over
q^2\alpha_0\right)^m\left(\sum_{j=1}^m{1\over j}\right)x^{qm}
$$
are solutions of
$$
\alpha_0x^2y''+\beta_0xy'+(\gamma_0+\gamma_qx^q)y=0.
$$
\end{alist}

\item\label{exer:7.6.56}
The equation
$$
xy''+y'+xy=0
$$
is
\href{https://mathshistory.st-andrews.ac.uk/Biographies/Bessel/}
 {\color{blue}\it Bessel's equation of order zero.\/} (See
Exercise~\ref{exer:7.5.53}.)
  Find two linearly independent Frobenius solutions of this equation.

\item\label{exer:7.6.57}
Suppose the assumptions of Exercise~7.5.~\hspace*{-3pt}\ref{exer:7.5.53}
hold, except that
$$
p_0(r)=\alpha_0(r-r_1)^2.
$$
 Show that
$$
y_1={x^{r_1}\over\alpha_0+\alpha_1x+\alpha_2x^2}\quad\mbox{ and }\quad
y_2={x^{r_1}\ln x\over\alpha_0+\alpha_1x+\alpha_2x^2}
$$
are linearly independent Frobenius  solutions of
$$
x^2(\alpha_0+\alpha_1x+\alpha_2 x^2)y''+x(\beta_0+\beta_1x+\beta_2x^2)y'+
(\gamma_0+\gamma_1x+\gamma_2x^2)y=0
$$
on any interval $(0,\rho)$ on which
$\alpha_0+\alpha_1x+\alpha_2x^2$ has no zeros.

\exercisetext{In Exercises~\ref{exer:7.6.58}--\ref{exer:7.6.65} use the method
suggested by Exercise~\ref{exer:7.6.57} to find the general solution on
some interval $(0,\rho)$.}

\item\label{exer:7.6.58}
$4x^2(1+x)y''+8x^2y'+(1+x)y=0$

\item\label{exer:7.6.59}
$9x^2(3+x)y''+3x(3+7x)y'+(3+4x)y=0$

\item\label{exer:7.6.60}
$x^2(2-x^2)y''-x(2+3x^2)y'+(2-x^2)y=0$

\item\label{exer:7.6.61}
$16x^2(1+x^2)y''+8x(1+9x^2)y'+(1+49x^2)y=0$

\item\label{exer:7.6.62}
$x^2(4+3x)y''-x(4-3x)y'+4y=0$

\item\label{exer:7.6.63}
$4x^2(1+3x+x^2)y''+8x^2(3+2x)y'+(1+3x+9x^2)y=0$

\item\label{exer:7.6.64}
$x^2(1-x)^2y''-x(1+2x-3x^2)y'+(1+x^2)y=0$

\item\label{exer:7.6.65}
$9x^2(1+x+x^2)y''+3x(1+7x+13x^2)y'+(1+4x+25x^2)y=0$

\item\label{exer:7.6.66}
\begin{alist}
\item % (a)
Let $L$ and $y(x,r)$ be as in Exercises~\ref{exer:7.5.57} and
\ref{exer:7.5.58}. Extend Theorem~\ref{thmtype:7.6.1} by showing that
$$
L\left({\partial y\over \partial
r}(x,r)\right)=p'_0(r)x^r+x^rp_0(r)\ln x.
$$
\item % (b)
Show that if
$$
p_0(r)=\alpha_0(r-r_1)^2
$$
then
$$
y_1=y(x,r_1) \mbox{\quad and \quad} y_2={\partial
y\over\partial r}(x,r_1)
$$
 are solutions of $Ly=0$.
\end{alist}

\end{exerciselist}
